% !TeX program = xelatex
% Associative memory: a math-intuition-first Chinese tutorial.
% Compile: latexmk -xelatex -interaction=nonstopmode -halt-on-error main.tex
\documentclass[11pt,a4paper,UTF8]{ctexart}
\usepackage[a4paper,top=22mm,bottom=23mm,left=23mm,right=23mm,headheight=22pt]{geometry}
\usepackage{fontspec}
\setCJKmainfont{Noto Serif CJK SC}
\setCJKsansfont{Noto Sans CJK SC}
\setCJKmonofont{Noto Sans CJK SC}
\setmainfont{Noto Serif}
\setsansfont{Noto Sans}
\usepackage{amsmath,amssymb,mathtools}
\usepackage{microtype}
\usepackage{booktabs,tabularx,array,multirow}
\usepackage{enumitem}
\usepackage{graphicx}
\usepackage{xcolor}
\usepackage{tikz}
\usetikzlibrary{arrows.meta,positioning,calc,fit,shapes.geometric,backgrounds}
\usepackage{pgfplots}
\pgfplotsset{compat=1.18}
\usepackage[most]{tcolorbox}
\usepackage{fancyhdr}
\usepackage{titlesec}
\usepackage{hyperref}
\usepackage{needspace}
\usepackage{float}

\definecolor{ink}{HTML}{172A40}
\definecolor{blue}{HTML}{147CA3}
\definecolor{teal}{HTML}{078D82}
\definecolor{orange}{HTML}{C27A31}
\definecolor{paper}{HTML}{F3F7FA}
\definecolor{linegray}{HTML}{D8E3EB}
\definecolor{muted}{HTML}{5B6F80}
\definecolor{blueback}{HTML}{EAF5F9}
\definecolor{orangeback}{HTML}{FCF5E9}
\definecolor{greenback}{HTML}{EAF7F3}
\color{ink}
\hypersetup{unicode=true,colorlinks=true,linkcolor=blue,urlcolor=blue,pdftitle={从矩阵乘法到 MLP 与 Attention：一条记忆的旅程},pdfauthor={学习讲义},pdfsubject={Associative Memory / Transformer 数学推导}}
\setlength{\parindent}{2em}
\setlength{\parskip}{2.7pt}
\linespread{1.19}
\setlist[itemize]{leftmargin=2.0em,itemsep=2pt,topsep=4pt}
\setlist[enumerate]{leftmargin=2.0em,itemsep=2pt,topsep=4pt}
\numberwithin{equation}{section}
\titleformat{\section}{\sffamily\Large\bfseries\color{ink}}{\color{blue}\thesection}{0.7em}{}
\titleformat{\subsection}{\sffamily\large\bfseries\color{ink}}{\color{blue}\thesubsection}{0.7em}{}
\titlespacing*{\section}{0pt}{15pt}{7pt}
\titlespacing*{\subsection}{0pt}{12pt}{6pt}
\pagestyle{fancy}
\fancyhf{}
\fancyhead[L]{\sffamily\footnotesize\color{muted} AI by Intuition · Chapter 01}
\fancyhead[R]{\sffamily\footnotesize\color{muted}\nouppercase{\leftmark}}
\fancyfoot[C]{\sffamily\footnotesize\color{muted}\thepage}
\renewcommand{\headrulewidth}{0.35pt}
\renewcommand{\headrule}{\hbox to\headwidth{\color{linegray}\leaders\hrule height \headrulewidth\hfill}}
\ctexset{section={name={,、},number=\chinese{section}}}

\newtcolorbox{idea}[1][]{enhanced,breakable,colback=blueback,colframe=blueback,boxrule=0pt,arc=2mm,left=4mm,right=4mm,top=3mm,bottom=3mm,borderline west={2.5pt}{0pt}{blue},fonttitle=\sffamily\bfseries,coltitle=ink,title={数学直觉},#1}
\newtcolorbox{example}[1][]{enhanced,breakable,colback=orangeback,colframe=orangeback,boxrule=0pt,arc=2mm,left=4mm,right=4mm,top=3mm,bottom=3mm,borderline west={2.5pt}{0pt}{orange},fonttitle=\sffamily\bfseries,coltitle=ink,title={具体例子},#1}
\newtcolorbox{conclusion}[1][]{enhanced,breakable,colback=greenback,colframe=greenback,boxrule=0pt,arc=2mm,left=4mm,right=4mm,top=3mm,bottom=3mm,borderline west={2.5pt}{0pt}{teal},fonttitle=\sffamily\bfseries,coltitle=ink,title={这一节记住什么},#1}
\newtcolorbox{warningbox}[1][]{enhanced,breakable,colback=paper,colframe=linegray,boxrule=.6pt,arc=1.5mm,left=3.5mm,right=3.5mm,top=2.5mm,bottom=2.5mm,#1}
\newcommand{\vect}[1]{\boldsymbol{#1}}
\newcommand{\relu}{\operatorname{ReLU}}
\newcommand{\softmax}{\operatorname{softmax}}
\newcommand{\R}{\mathbb{R}}
\newcommand{\eyebrow}[1]{{\sffamily\bfseries\footnotesize\color{blue}#1}}

\begin{document}
\thispagestyle{empty}
\vspace*{19mm}
\noindent\colorbox{blue}{\parbox{47mm}{\centering\color{white}\sffamily\bfseries AI by Intuition · 01}}\par
\vspace{11mm}
{\sffamily\bfseries\fontsize{26}{34}\selectfont 为什么矩阵乘法\\[4mm]可以检索记忆？}\par
\vspace{7mm}
{\sffamily\fontsize{14}{21}\selectfont 从二维点积、Key--Value 堆叠出发，一步步“重新发明” Transformer}\par
\vspace{5mm}
\noindent{\color{muted}\rule{\linewidth}{0.7pt}}\par
\vspace{8mm}
\begin{tcolorbox}[enhanced,colback=paper,colframe=paper,arc=3mm,boxrule=0pt,left=6mm,right=6mm,top=6mm,bottom=6mm]
\textbf{只需要这些预备知识：}加减乘除、平方与平方根、基本的勾股定理。\par
\textbf{学习路线：} 向量 $\rightarrow$ 点积 $\rightarrow$ 外积 $\rightarrow$ Key/Value 堆叠 $\rightarrow$ ReLU 与 MLP $\rightarrow$ Attention $\rightarrow$ 记忆容量。\par
\textbf{目标：} 看懂公式并知道\emph{为什么需要它}，而不是背诵符号。
\end{tcolorbox}
\vspace{9mm}
\begin{center}
\begin{tikzpicture}[>=Latex, node distance=1.15cm and 0.6cm,
    n/.style={draw=linegray,fill=white,rounded corners=4pt,minimum width=35mm,minimum height=13mm,align=center,font=\sffamily\bfseries}]
\node[n] (lin) {Linear\\[-1mm] $\vect{y}=W\vect{x}$};
\node[n,right=of lin] (mlp) {MLP\\[-1mm] $\vect{y}=V^\top\sigma(K\vect{x})$};
\node[n,right=of mlp] (att) {Attention\\[-1mm] $\vect{y}=V^\top\softmax(\beta K\vect{x})$};
\draw[->,thick,blue] (lin)--(mlp) node[midway,above,font=\footnotesize,align=center]{压低干扰};
\draw[->,thick,blue] (mlp)--(att) node[midway,above,font=\footnotesize,align=center]{选择更强};
\end{tikzpicture}
\end{center}
\vfill
\noindent\textbf{参考文章}\quad \textit{You could have invented transformers from memory}\par
\noindent\href{https://alexwang.ai/posts/associative-memory}{https://alexwang.ai/posts/associative-memory}\par
\vspace{3mm}
\noindent {\color{muted}\small AI by Intuition 第 01 篇 · 先有问题，再有直觉，然后才有公式。}\par
\clearpage
\tableofcontents
\clearpage

\section{向量：如何用数字表示一个东西？}

假设我们要让计算机记住两条关联：\textbf{法国 $\rightarrow$ 巴黎}，\textbf{日本 $\rightarrow$ 东京}。人类看得懂文字，但公式需要能计算的对象，所以先给每个“检索线索”一个数字表示。

例如，将法国和日本分别编码为：
\begin{equation}
\vect{k}_{\text{法国}}=\begin{bmatrix}1\\0\end{bmatrix},\qquad
\vect{k}_{\text{日本}}=\begin{bmatrix}0\\1\end{bmatrix}.
\end{equation}

这种有序数字列表叫作\textbf{向量（vector）}。$(1,0)$ 是二维向量；真实模型中的向量可以有几十、几百乃至几千维。这里的数字只是教学示意，并非语言模型真实的语义编码。

\begin{center}
\begin{tikzpicture}[scale=1.02,>=Latex]
\draw[->,gray!70] (-.2,0)--(5.0,0) node[below] {$x$};
\draw[->,gray!70] (0,-.2)--(0,3.0) node[left] {$y$};
\draw[->,line width=1.6pt,blue] (0,0)--(4.0,0) node[above right] {法国 $(1,0)$};
\draw[->,line width=1.6pt,orange] (0,0)--(0,2.45) node[right] {日本 $(0,1)$};
\fill[ink] (0,0) circle(1.2pt);\node[below left] at (0,0) {$0$};
\end{tikzpicture}
\end{center}

向量还有一个重要性质：\textbf{长度}。二维向量 $\vect{a}=(a_1,a_2)$ 的长度就是勾股定理：
\begin{equation}
\|\vect{a}\|=\sqrt{a_1^2+a_2^2}.
\end{equation}
例如 $(0.6,0.8)$ 的长度为 $\sqrt{0.6^2+0.8^2}=1$。长度等于 $1$ 的向量叫\textbf{单位向量}，后面很多计算会方便得多。

\begin{idea}
不要把向量想得太复杂。它就是一个可以进行几何和代数运算的\textbf{数字地址}。我们稍后用它代表记忆的 Key（检索线索），用另一个向量代表 Value（记忆内容）。
\end{idea}

\section{点积：如何判断两条线索相不相似？}

给定 $\vect{a}=(a_1,a_2)$ 和 $\vect{b}=(b_1,b_2)$，它们的\textbf{点积}定义为：
\begin{equation}
\boxed{\vect{a}\cdot\vect{b}=a_1b_1+a_2b_2.}
\end{equation}
只要把对应数字相乘再加起来即可。例如：
\[
(1,0)\cdot(0.6,0.8)=1\times0.6+0\times0.8=0.6.
\]

\subsection{为什么点积代表相似度？}
点积还可以写成：
\begin{equation}
\vect{a}\cdot\vect{b}=\|\vect{a}\|\,\|\vect{b}\|\cos\theta,
\end{equation}
$\theta$ 是两个箭头的夹角。如果它们都是单位向量，就有：
\begin{equation}
\boxed{\vect{a}\cdot\vect{b}=\cos\theta.}
\end{equation}
方向越一致，点积越接近 $1$；垂直时点积等于 $0$；完全反方向时点积等于 $-1$。

\begin{center}
\begin{tikzpicture}[>=Latex]
\foreach \x/\ang/\score/\lab in {0/17/{\approx 1}/同向附近,4.2/90/{0}/垂直,8.4/165/{\approx -1}/反方向附近} {
 \begin{scope}[shift={(\x,0)}]
 \draw[->,thick,blue] (0,0)--(2.2,0);
 \draw[->,thick,orange] (0,0)--({2.1*cos(\ang)},{2.1*sin(\ang)});
 \fill (0,0) circle(1pt);
 \node[font=\footnotesize,align=center] at (1.05,-0.55) {\lab\\点积 $\score$};
 \end{scope}
}
\end{tikzpicture}
\end{center}

设某条记忆的 Key 是 $\vect{k}_i$，当前查询是 $\vect{x}$。最简单的匹配分数就是：
\begin{equation}
\boxed{s_i=\vect{k}_i^\top\vect{x}.}
\end{equation}
这里上标 $\top$ 表示转置；$\vect{k}_i^\top\vect{x}$ 与 $\vect{k}_i\cdot\vect{x}$ 都是点积。

\begin{conclusion}
点积就是最基础的\textbf{“匹配多少”}。Attention 中的 $QK^\top$，本质上也是在批量计算 Query 和各个 Key 的匹配程度。
\end{conclusion}

\section{外积与矩阵：如何真正“写入”一条记忆？}

现在为一条记忆指定两个向量：
\[
\text{Key }\vect{k}=\begin{bmatrix}1\\0\end{bmatrix},\qquad
\text{Value }\vect{v}=\begin{bmatrix}2\\3\end{bmatrix}.
\]
我们的目标是设计函数 $f$，使输入 Key 时返回 Value，即 $f(\vect{k})=\vect{v}$。

\subsection{先造出一个能工作的公式}
考虑：
\begin{equation}
\boxed{f(\vect{x})=\vect{v}\bigl(\vect{k}^\top\vect{x}\bigr).}
\end{equation}
如果输入 $\vect{x}=\vect{k}$，那么两单位向量的点积等于 $1$，所以：
\[
f(\vect{k})=\vect{v}\times 1=\vect{v}.
\]

\begin{center}
\begin{tikzpicture}[>=Latex,node distance=1cm,n/.style={draw=linegray,fill=paper,rounded corners=4pt,align=center,minimum height=13mm,minimum width=34mm,font=\small}]
\node[n] (a) {输入 $\vect{x}$};
\node[n,right=of a] (b) {$\vect{k}^\top\vect{x}$\\算匹配分数};
\node[n,right=of b] (c) {$\vect{v}(\vect{k}^\top\vect{x})$\\得到输出};
\draw[->,thick,blue] (a)--(b);\draw[->,thick,blue] (b)--(c);
\end{tikzpicture}
\end{center}

这个计算已经很像记忆检索：先看输入和哪条 Key 匹配，再按照匹配程度返回 Value。

\subsection{把关系装进矩阵：外积}
两向量的\textbf{外积（outer product）}是：
\begin{equation}
\vect{v}\vect{k}^\top=
\begin{bmatrix}2\\3\end{bmatrix}
\begin{bmatrix}1&0\end{bmatrix}
=\begin{bmatrix}2&0\\3&0\end{bmatrix}.
\end{equation}
外积与点积的区别：点积产生\textbf{一个数}；外积产生\textbf{一个矩阵}。

记这个矩阵为 $W=\vect{v}\vect{k}^\top$。根据矩阵乘法结合律：
\begin{equation}
\boxed{f(\vect{x})=\vect{v}(\vect{k}^\top\vect{x})=(\vect{v}\vect{k}^\top)\vect{x}=W\vect{x}.}
\end{equation}
这说明：我们可以把一条 Key--Value 关联“写入”权重矩阵 $W$，再用普通矩阵乘法取回来。

\begin{idea}
\textbf{一个线性层 $W\vect{x}$，可以被看作一个记忆检索器。}这正是文章把神经网络的线性层与联想记忆联系起来的第一步。
\end{idea}

\section{多条记忆：为什么会互相干扰？}

如果有 $M$ 条 Key--Value 关系 $(\vect{k}_i,\vect{v}_i)$，最简单的存储办法是\textbf{把外积全部相加}：
\begin{equation}
\boxed{W=\sum_{i=1}^{M}\vect{v}_i\vect{k}_i^\top.}
\end{equation}
查询时：
\begin{equation}
\boxed{W\vect{x}=\sum_{i=1}^{M}\vect{v}_i(\vect{k}_i^\top\vect{x}).}
\end{equation}
也就是说：\textbf{所有记忆都会按照它们与输入的相似度，贡献一部分输出。}

\subsection{一个完全算得出来的二维例子}
假设我们有两条记忆：
\begin{align*}
\vect{k}_A&=\begin{bmatrix}1\\0\end{bmatrix}, &
\vect{v}_A&=\begin{bmatrix}1\\0\end{bmatrix},\\
\vect{k}_B&=\begin{bmatrix}0.6\\0.8\end{bmatrix}, &
\vect{v}_B&=\begin{bmatrix}0\\1\end{bmatrix}.
\end{align*}
两条记忆分别形成一个矩阵：
\[
\vect{v}_A\vect{k}_A^\top=
\begin{bmatrix}1&0\\0&0\end{bmatrix},\qquad
\vect{v}_B\vect{k}_B^\top=
\begin{bmatrix}0&0\\0.6&0.8\end{bmatrix}.
\]
把它们加起来：
\begin{equation}
\boxed{W=\begin{bmatrix}1&0\\0.6&0.8\end{bmatrix}.}
\end{equation}

现在查询 A，即 $\vect{x}=\vect{k}_A$：
\begin{equation}
W\vect{k}_A=
\begin{bmatrix}1&0\\0.6&0.8\end{bmatrix}
\begin{bmatrix}1\\0\end{bmatrix}
=\begin{bmatrix}1\\0.6\end{bmatrix}.
\end{equation}
然而我们真正想得到的结果是 $\vect{v}_A=(1,0)^\top$。

\begin{example}
计算结果分成两部分：
\[
W\vect{k}_A=\underbrace{\begin{bmatrix}1\\0\end{bmatrix}}_{\text{正确记忆}}
+\underbrace{\begin{bmatrix}0\\0.6\end{bmatrix}}_{\text{来自 B 的干扰}}.
\]
为何有 $0.6$？因为 $\vect{k}_B^\top\vect{k}_A=0.6$，即 B 的 Key 与查询 A 并不垂直，所以 B 也被激活了。
\end{example}

\begin{center}
\begin{tikzpicture}
\node[anchor=east,font=\small] at (0,0.65) {A（正确信号）};
\node[anchor=east,font=\small] at (0,-.07) {B（干扰）};
\fill[paper] (0,0.45) rectangle (7.5,0.8);
\fill[paper] (0,-.25) rectangle (7.5,0.1);
\fill[blue] (0,0.45) rectangle (7.5,0.8);
\fill[orange] (0,-.25) rectangle (4.5,0.1);
\node[anchor=west,font=\small] at (7.7,.62) {$1.0$};
\node[anchor=west,font=\small] at (4.7,-.07) {$0.6$};
\end{tikzpicture}
\end{center}

对第 $j$ 条单位 Key 进行查询时，通常可以写为：
\begin{equation}
W\vect{k}_j=
\underbrace{\vect{v}_j}_{\text{正确信号}}
+\underbrace{\sum_{i\ne j}\vect{v}_i(\vect{k}_i^\top\vect{k}_j)}_{\text{其余记忆的干扰}}.
\end{equation}

\begin{conclusion}
关键矛盾不是\textbf{能不能把记忆写进矩阵}，而是\textbf{取回其中一条时，其他记忆会不会一起冒出来}。后面的 MLP 和 Attention 都可以理解成对这个问题的不同改进。
\end{conclusion}

\section{改进一：让线性记忆写得更聪明}

我们刚才直接将外积相加，写入新记忆时没有考虑旧记忆。其实也可以寻找一个更好的矩阵 $W$，要求它尽可能满足所有配对关系：
\begin{equation}
\boxed{\min_W\sum_{i=1}^{M}\|W\vect{k}_i-\vect{v}_i\|^2.}
\end{equation}
这叫\textbf{最小二乘}：选择一个矩阵，使总平方误差尽可能小。

在上面的二维例子中，外积相加给出：
\[
W_{\text{simple}}=\begin{bmatrix}1&0\\0.6&0.8\end{bmatrix}.
\]
但我们也可以找到：
\begin{equation}
W_{\text{better}}=\begin{bmatrix}1&-0.75\\0&1.25\end{bmatrix}.
\end{equation}
直接代入就可以验证：
\[
W_{\text{better}}\vect{k}_A=\vect{v}_A,\qquad
W_{\text{better}}\vect{k}_B=\vect{v}_B.
\]
两条记忆都可以精确召回！矩阵中的负数在这里扮演了\textbf{抵消交叉干扰}的角色。

不过，线性映射的自由度有限。如果在 $d_k$ 维中放入很多\emph{任意且相互独立的} Key--Value 对，最终还是无法全部精确拟合。

\begin{idea}
提高记忆能力有两种思路：\textbf{写入时消除干扰}（优化矩阵），以及\textbf{检索时抑制干扰}（设计更好的权重函数）。接下来采用第二种。
\end{idea}


\section{把 Key 和 Value 堆成矩阵：中间并没有魔法}

我们已经知道每条记忆都有一个 Key 和一个 Value。但初学者最容易卡住的恰恰是
\[
\sum_i\vect{v}_i\,\sigma(\vect{k}_i^\top\vect{x})
\quad\longrightarrow\quad
V^\top\sigma(K\vect{x}).
\]
\textbf{为什么能堆？为什么 ReLU 能放在两个矩阵乘法之间？}这一节完全拆开手算。

\subsection{两条记忆，各用两个数字表示}

先约定两条 Key--Value 关系：
\begin{align*}
\vect{k}_1&=\begin{bmatrix}1\\0\end{bmatrix},
& \vect{v}_1&=\begin{bmatrix}2\\1\end{bmatrix},\\
\vect{k}_2&=\begin{bmatrix}0.6\\0.8\end{bmatrix},
& \vect{v}_2&=\begin{bmatrix}0\\3\end{bmatrix}.
\end{align*}
给出查询：
\[
\vect{x}=\begin{bmatrix}1\\-1\end{bmatrix}.
\]
\textbf{不使用矩阵}时，分别计算两个匹配分数：
\begin{align*}
s_1&=\vect{k}_1^\top\vect{x}=1\cdot1+0\cdot(-1)=1,\\
s_2&=\vect{k}_2^\top\vect{x}=0.6\cdot1+0.8\cdot(-1)=-0.2.
\end{align*}
所以两个 Key 的匹配程度为 $1$ 和 $-0.2$。现在我们只想把\textbf{这两次点积合并计算}。

\subsection{第一个矩阵：每行是一条 Key}

把两条 Key \textbf{横放成两行}：
\begin{equation}
K=\begin{bmatrix}
  1 & 0\\
  0.6 & 0.8
\end{bmatrix}.
\end{equation}
矩阵乘列向量的规则是：\textbf{左边每一行，与右边那一列分别点积}。

\begin{example}[title={用一行乘一列，亲手算两次}]
\begin{align*}
K\vect{x}
&=\begin{bmatrix}1&0\\0.6&0.8\end{bmatrix}
  \begin{bmatrix}1\\-1\end{bmatrix}\\
&=\begin{bmatrix}
1\cdot1+0\cdot(-1)\\
0.6\cdot1+0.8\cdot(-1)
\end{bmatrix}
=\boxed{\begin{bmatrix}1\\-0.2\end{bmatrix}}.
\end{align*}
第一行的结果恰好是 $s_1$，第二行的结果恰好是 $s_2$。
\end{example}

因此，下面两个表达式\textbf{一模一样}：
\[
K\vect{x}
=
\begin{bmatrix}\vect{k}_1^\top\vect{x}\\\vect{k}_2^\top\vect{x}\end{bmatrix}
=
\begin{bmatrix}s_1\\s_2\end{bmatrix}.
\]
堆成矩阵并没有添加新的数学假设，只是用一种方便批量计算的记法表示多次点积。

\begin{center}
\begin{tikzpicture}[>=Latex,
 box/.style={draw=linegray,rounded corners=3pt,fill=paper,minimum height=11mm,align=center,font=\small},
 node distance=8mm]
\node[box,minimum width=38mm] (key) {$K$：两行 Key};
\node[box,right=of key,minimum width=33mm] (query) {$\vect{x}$：一列 Query};
\node[box,right=of query,minimum width=41mm] (score) {$K\vect{x}$：两项分数};
\draw[->,very thick,blue] (key)--(query);
\draw[->,very thick,blue] (query)--(score);
\end{tikzpicture}
\end{center}

\subsection{ReLU 不乘矩阵：它逐个处理分数}

矩阵乘法算完后，结果已经是一个普通的两维向量：
\[
\vect{s}=K\vect{x}=\begin{bmatrix}1\\-0.2\end{bmatrix}.
\]
激活函数 $\relu(t)=\max(0,t)$ 只处理\textbf{一个数}。当输入是向量时，约定它\textbf{逐元素处理}：
\begin{equation}
\relu(\vect{s})
=\begin{bmatrix}\relu(1)\\\relu(-0.2)\end{bmatrix}
=\boxed{\begin{bmatrix}1\\0\end{bmatrix}}.
\end{equation}
这一步不是矩阵乘法，也不是对 $K$ 的每个元素先做 ReLU。\textbf{它处理的是 $K\vect{x}$ 计算出的两项匹配分数。}

记新的分数向量为
\[
\vect{a}=\relu(K\vect{x})=\begin{bmatrix}1\\0\end{bmatrix}.
\]
直觉是：第一条记忆保留权重 $1$；第二条记忆被压到 $0$。

\subsection{第二个矩阵：每列是一条 Value}

我们希望的输出是\textbf{按照激活分数，将 Value 加权求和}：
\begin{equation}
\vect{y}=a_1\vect{v}_1+a_2\vect{v}_2.
\end{equation}
现在把两个 Value 排成\textbf{两列}：
\begin{equation}
V^\top=
\begin{bmatrix}
\vect{v}_1&\vect{v}_2
\end{bmatrix}
=
\begin{bmatrix}
2&0\\
1&3
\end{bmatrix}.
\end{equation}
为什么是列？因为矩阵乘列向量也有等价的解释：\textbf{用右侧的每个数字，对左侧相应列做加权求和}。

\begin{example}[title={矩阵乘向量 = 按列加权求和}]
\begin{align*}
V^\top\vect{a}
&=\begin{bmatrix}2&0\\1&3\end{bmatrix}
  \begin{bmatrix}1\\0\end{bmatrix}\\
&=\begin{bmatrix}2\cdot1+0\cdot0\\1\cdot1+3\cdot0\end{bmatrix}\\
&=1\begin{bmatrix}2\\1\end{bmatrix}
 +0\begin{bmatrix}0\\3\end{bmatrix}
=\boxed{\begin{bmatrix}2\\1\end{bmatrix}}.
\end{align*}
我们取回了第一条 Value，第二条因为激活为 $0$ 而没有参与。
\end{example}

\subsection{整个过程连在一起}

现在把刚才三个独立步骤写在同一张纸上：
\begin{align}
\text{匹配：}\quad
K\vect{x}
&=\begin{bmatrix}1\\-0.2\end{bmatrix},\\
\text{筛选：}\quad
\relu(K\vect{x})
&=\begin{bmatrix}1\\0\end{bmatrix},\\
\text{读出：}\quad
V^\top\relu(K\vect{x})
&=\begin{bmatrix}2\\1\end{bmatrix}.
\end{align}

\begin{center}
\begin{tikzpicture}[>=Latex,node distance=5mm,
 n/.style={draw=linegray,rounded corners=3pt,fill=paper,minimum width=102mm,minimum height=11mm,align=center,font=\small}]
\node[n] (a) {输入 $\vect{x}=(1,-1)^\top$};
\node[n,below=of a] (b) {$K\vect{x}=(1,-0.2)^\top$：对所有 Key 算匹配分数};
\node[n,below=of b] (c) {$\relu(K\vect{x})=(1,0)^\top$：逐元素过滤};
\node[n,below=of c] (d) {$V^\top\relu(K\vect{x})=(2,1)^\top$：按列混合 Value};
\draw[->,thick,blue] (a)--(b);
\draw[->,thick,blue] (b)--(c);
\draw[->,thick,blue] (c)--(d);
\end{tikzpicture}
\end{center}

\begin{conclusion}
一定要分开理解：\textbf{$K\vect{x}$ 是批量点积；ReLU 是逐元素函数；$V^\top\vect{a}$ 是按列加权求和。}每一步的输出都恰好是下一步能够接收的输入。
\end{conclusion}

\subsection{维度检查：为什么它们可以相乘？}

前面两条记忆、二维 Key、二维 Value，所以恰好都是 $2\times2$ 方阵。但真实情况并不需要是方阵。

假设共有 $M$ 条记忆；每个 Key 有 $d_k$ 个数字，每个 Value 有 $d_v$ 个数字。我们约定：
\[
\underbrace{K}_{M\times d_k}
\underbrace{\vect{x}}_{d_k\times1}
=\underbrace{\vect{s}}_{M\times1},
\qquad
\underbrace{V^\top}_{d_v\times M}
\underbrace{\sigma(\vect{s})}_{M\times1}
=\underbrace{\vect{y}}_{d_v\times1}.
\]
从矩阵大小看，\textbf{内侧维度必须相等}：第一个乘法的 $d_k$ 对得上；第二个乘法的 $M$ 对得上。

例如 100 条记忆，Key 维度 64，Value 维度 32：
\[
(100\times64)(64\times1)=(100\times1),
\]
\[
(32\times100)(100\times1)=(32\times1).
\]
这里并没有神秘的“特殊神经网络矩阵乘法”，始终是普通的线性代数。

\subsection{这和外积求和到底是什么关系？}

如果\textbf{不加 ReLU}，之前的外积求和能合并：
\begin{equation}
\underbrace{V^\top K}_{W}
=
\begin{bmatrix}\vect{v}_1&\vect{v}_2\end{bmatrix}
\begin{bmatrix}\vect{k}_1^\top\\\vect{k}_2^\top\end{bmatrix}
=
\vect{v}_1\vect{k}_1^\top+\vect{v}_2\vect{k}_2^\top.
\end{equation}
所以线性召回可以写成：
\[
W\vect{x}=(V^\top K)\vect{x}=V^\top(K\vect{x}).
\]
\textbf{一旦把 ReLU 放在中间，通常就不能合并成一个固定的 $W$。}以本节数字为例：
\begin{align*}
V^\top K\vect{x}
&=1\vect{v}_1-0.2\vect{v}_2
=\begin{bmatrix}2\\0.4\end{bmatrix},\\
V^\top\relu(K\vect{x})
&=1\vect{v}_1+0\vect{v}_2
=\begin{bmatrix}2\\1\end{bmatrix}.
\end{align*}
前者保留了负匹配；后者把它滤掉了。这就是激活函数引入\textbf{非线性}的实际效果。

\begin{warningbox}
\textbf{约定提醒：}本讲义规定 $K$ 与 $V$ 的每一\textbf{行}是一条 Key、Value，所以公式使用 $V^\top\sigma(K\vect{x})$。有些教材直接把 Value 放在\textbf{列}中，记为 $\widetilde V$，便会写成 $\widetilde V\sigma(K\vect{x})$。只是摆放约定不同，内容相同。标准 Transformer 的批量 $Q,K,V$ 通常按 Token 放在行上。
\end{warningbox}

\section{改进二：非线性为什么会导向 MLP？}

之前每条记忆的贡献是：
\[
\vect{v}_i\times\underbrace{(\vect{k}_i^\top\vect{x})}_{\text{原始相似度}}.
\]
例如正确 Key 的相似度是 $1$，无关 Key 是 $0.6$。线性加权得到 $1$ 和 $0.6$，错误记忆仍贡献很多。

最简单的办法：\textbf{把相似度平方}！
\[
1^2=1,\qquad 0.6^2=0.36.
\]
如果立方：
\[
1^3=1,\qquad 0.6^3=0.216.
\]
正确记忆保持为 $1$，干扰越来越小。

\begin{center}
\begin{tikzpicture}
\begin{axis}[
 width=.89\linewidth,height=5.4cm,
 axis lines=left,xlabel={相似度 $s$},ylabel={激活权重},
 xmin=0,xmax=1.02,ymin=0,ymax=1.06,
 xtick={0,.2,.4,.6,.8,1},ytick={0,.2,.4,.6,.8,1},
 legend style={at={(.02,.98)},anchor=north west,font=\small,draw=none,fill=white},
 grid=major,grid style={linegray!55}, tick label style={font=\small},
 ]
 \addplot[blue,very thick,domain=0:1,samples=100] {x};
 \addlegendentry{线性 $s$}
 \addplot[orange,very thick,domain=0:1,samples=100] {x^2};
 \addlegendentry{平方 $s^2$}
 \addplot[teal,very thick,domain=0:1,samples=100] {x^3};
 \addlegendentry{立方 $s^3$}
 \addplot[dashed,gray] coordinates {(.6,0)(.6,1)};
\end{axis}
\end{tikzpicture}
\end{center}

为了把负相似度也压成零，可以使用：
\begin{equation}
\sigma(s)=\relu(s)^2=\bigl(\max(0,s)\bigr)^2.
\end{equation}
我们于是得到新的记忆检索器：
\begin{equation}
\boxed{f(\vect{x})=\sum_i\vect{v}_i\,\sigma(\vect{k}_i^\top\vect{x}).}
\end{equation}
将所有 Key 堆成矩阵 $K$，所有 Value 堆成矩阵 $V$，就可以简写为：
\begin{equation}
\boxed{f(\vect{x})=V^\top\sigma(K\vect{x}).}
\end{equation}
这里约定 $K$ 的每一行是一条 Key，$V$ 的每一行是一条 Value；激活函数 $\sigma$ 对向量中的各元素分别作用。

\begin{center}
\begin{tikzpicture}[>=Latex,node distance=0.55cm,n/.style={draw=linegray,rounded corners=3pt,fill=paper,minimum width=92mm,minimum height=11mm,align=center,font=\small}]
\node[n] (q) {输入 Query $\vect{x}$};
\node[n,below=of q] (s) {第一层：$K\vect{x}$，计算每条记忆的匹配分数};
\node[n,below=of s] (a) {非线性：$\sigma(K\vect{x})$，压低不相关记忆的激活};
\node[n,below=of a] (v) {第二层：$V^\top\sigma(K\vect{x})$，组合 Value 得到输出};
\draw[->,blue,thick] (q)--(s);\draw[->,blue,thick] (s)--(a);\draw[->,blue,thick] (a)--(v);
\end{tikzpicture}
\end{center}

这恰好是一个\textbf{两层 MLP} 的典型数学形式：一次线性变换、一次非线性激活、再一次线性变换。

\begin{conclusion}
从记忆角度看，MLP 的第一层可以被解释为\textbf{匹配模式}，激活函数负责\textbf{筛选}，第二层负责\textbf{读出内容}。这是一种有启发性的解释，并不意味着每个真实神经元都恰好对应一条可分离的事实。
\end{conclusion}

\section{改进三：从指数函数走到 Softmax Attention}

平方和立方都可以抑制弱匹配。如果想要更强的区分能力，可以考虑\textbf{指数函数}：
\begin{equation}
\sigma(s)=e^{\beta s},\qquad e\approx2.718.
\end{equation}
其中 $\beta$ 是“锐度”参数：越大，较高相似度相对于较低相似度的优势就越明显。

\subsection{把指数分数变成权重}
假设两条记忆的相似度仍然是：
\[
s_A=1,\qquad s_B=0.6.
\]
指数变换得到 $e^{\beta}$ 和 $e^{0.6\beta}$。为了让两个权重的总和等于 $1$，将每个指数值除以指数值之和：
\begin{align}
p_A&=\frac{e^{\beta}}{e^{\beta}+e^{0.6\beta}},\\
p_B&=\frac{e^{0.6\beta}}{e^{\beta}+e^{0.6\beta}}.
\end{align}
这就是 \textbf{Softmax}。

\begin{table}[H]
\centering
\caption{调整 $\beta$ 时的 Softmax 权重（两个分数为 $1$ 和 $0.6$）}
\begin{tabular}{@{}rrr@{}}
\toprule
$\beta$ & $p_A$：正确记忆 & $p_B$：干扰记忆\\
\midrule
$0$ & $50.0\%$ & $50.0\%$\\
$2$ & $69.0\%$ & $31.0\%$\\
$5$ & $88.1\%$ & $11.9\%$\\
$10$ & $98.2\%$ & $1.8\%$\\
$15$ & $99.8\%$ & $0.2\%$\\
\bottomrule
\end{tabular}
\end{table}

\begin{idea}
Softmax 没有“创造”新记忆，而是调整了\textbf{从已有记忆读取多少内容}。当 $\beta$ 很大时，权重越来越集中在相似度最高的 Key 上。
\end{idea}

\subsection{写出 Attention 的形式}
给第 $i$ 条记忆的权重为：
\begin{equation}
p_i=\frac{e^{\beta\vect{k}_i^\top\vect{x}}}{\sum_j e^{\beta\vect{k}_j^\top\vect{x}}}.
\end{equation}
最终输出就是各个 Value 的加权平均：
\begin{equation}
\boxed{f(\vect{x})=\sum_i p_i\vect{v}_i
=\sum_i\vect{v}_i\frac{e^{\beta\vect{k}_i^\top\vect{x}}}{\sum_j e^{\beta\vect{k}_j^\top\vect{x}}}.}
\end{equation}
矩阵写法：
\begin{equation}
\boxed{f(\vect{x})=V^\top\softmax(\beta K\vect{x}).}
\end{equation}
这就是\textbf{按相似度检索 Value} 的 Attention 核心公式。

在标准 Transformer 中，对一批 Query 同时计算，并使用广为人知的形式：
\begin{equation}
\boxed{\operatorname{Attention}(Q,K,V)=\softmax\!\left(\frac{QK^\top}{\sqrt{d_k}}\right)V.}
\end{equation}
三个字母的含义非常直观：
\begin{itemize}
\item $Q$，Query：我正在找什么？
\item $K$，Key：每条候选信息用什么线索被找到？
\item $V$，Value：被找到以后实际取回什么内容？
\end{itemize}
标准 Transformer 中用 $1/\sqrt{d_k}$ 缩放\emph{未单位归一化}的点积；前面用 $\beta$ 的例子多在\emph{单位向量}的设定下讨论。两者不能在不同的向量尺度之间直接对比数字大小。

\begin{center}
\begin{tikzpicture}[>=Latex, node distance=5mm,
 block/.style={draw=linegray,rounded corners=3pt,fill=paper,align=center,minimum height=10mm,font=\small}]
\node[block,minimum width=30mm] (q) {Query $Q$};
\node[block,right=17mm of q,minimum width=30mm] (k) {Keys $K$};
\node[block,right=17mm of k,minimum width=30mm] (v) {Values $V$};
\node[block,below=13mm of k,minimum width=40mm] (dot) {相似度 $QK^\top$};
\node[block,below=of dot,minimum width=40mm] (soft) {Softmax 权重};
\node[block,below=of soft,minimum width=40mm] (out) {加权求和：输出};
\draw[->,blue,thick] (q.south) |- (dot.west);
\draw[->,blue,thick] (k.south)--(dot.north);
\draw[->,blue,thick] (dot)--(soft);
\draw[->,blue,thick] (soft)--(out);
\draw[->,orange,thick] (v.south) |- (out.east);
\end{tikzpicture}
\end{center}

\begin{warningbox}
\textbf{一个重要区别：}MLP 的 Key/Value 视角通常来自\textbf{训练好的参数}；Self-Attention 的 Key/Value 则是由\textbf{当前上下文的 Token 表示动态产生}。所以两者虽然有相似的检索数学，但保存和更新信息的方式不同。
\end{warningbox}

\section{为什么维度影响记忆容量？}

接下来理解文章中看起来最神奇的数字：为什么一个 $64$ 维 Key，理论上可能在某些检索模型里对应远超 $64$ 条的记忆？

\subsection{随机高维向量往往近似正交}
如果两个独立随机单位向量在 $d_k$ 维空间中，那么它们点积的典型尺度为：
\begin{equation}
|\vect{k}_i^\top\vect{k}_j|\sim\frac{1}{\sqrt{d_k}},\qquad i\ne j.
\end{equation}
这个 $\sim$ 表示量级关系，不是每一对随机向量都恰好等于该值。

当 $d_k=64$，典型干扰相似度量级为：
\[
\frac{1}{\sqrt{64}}=\frac18=0.125.
\]
空间维度越高，随机无关 Key 往往越接近垂直，单条记忆造成的干扰通常越小。

\subsection{为什么有 \texorpdfstring{$\sqrt M$}{√M}？}
想象随机向左或向右走，每一步长度为 $1$。走 $100$ 步，并不是必然离原点 $100$ 步；方向会互相抵消，最终位移的典型大小是 $\sqrt{100}=10$。

如果各条干扰的方向、符号足够随机且近似独立，那么 $M$ 个小干扰的总波动规模通常不是 $M$ 倍，而是 $\sqrt M$ 倍。

结合上一节：
\begin{equation}
\boxed{\frac{\text{干扰}}{\text{信号}}\ \sim\ \sqrt M\times\frac{1}{\sqrt{d_k}}
=\sqrt{\frac{M}{d_k}}.}
\end{equation}
当干扰与信号处于同一量级时：
\begin{equation}
\sqrt{\frac{M}{d_k}}\sim1\quad\Rightarrow\quad M\sim d_k.
\end{equation}
因此在这个\textbf{随机 Key/Value、简单外积求和}的模型里，容量的标度大约是维度 $d_k$。

\subsection{平方激活为何带来平方量级？}
对于不相关 Key，原始相似度的典型大小是 $1/\sqrt{d_k}$；平方后，它变成约 $1/d_k$。在一些独立、均值抵消等条件成立的理想模型中，随机干扰的总体量级可以粗略看成：
\[
\frac{\text{干扰}}{\text{信号}}\ \sim\ \frac{\sqrt M}{d_k}.
\]
设这个量级达到 $1$：
\begin{equation}
\frac{\sqrt M}{d_k}\sim1
\quad\Rightarrow\quad
\boxed{M\sim d_k^2.}
\end{equation}
这就解释了\textbf{为什么改变相似度的非线性，可以改变理论容量对维度的增长规律}。一般的 $p$ 次选择性激活在相应理想假设下可产生 $M\sim d_k^p$ 这一类标度。

\begin{warningbox}
\textbf{请特别注意：}这些只是文章讨论的模型下的容量标度和噪声启发式解释。真实 ReLU 平方激活、非零均值、相关 Key、非独立 Value、非精确 Key 查询等因素都会改变结论。\textbf{“记得住的向量关联数”也不等于“LLM 可靠知道的事实数”。}
\end{warningbox}

\section{回头看：Linear、MLP、Attention 究竟改变了什么？}

\begin{table}[H]
\centering
\renewcommand{\arraystretch}{1.45}
\caption{三类模型的统一检索视角}
\begin{tabularx}{\linewidth}{@{}>{\bfseries}p{21mm}X X@{}}
\toprule
模型 & 读出规则 & 直观意义 \\
\midrule
Linear & $\sum_i \vect{v}_i (\vect{k}_i^\top\vect{x})$ & 任何非零相似度都参与输出 \\
MLP & $\sum_i \vect{v}_i\sigma(\vect{k}_i^\top\vect{x})$ & 非线性抑制弱匹配 \\
Attention & $\sum_i \vect{v}_i\softmax(\beta K\vect{x})_i$ & 用归一化指数权重聚焦匹配项 \\
\bottomrule
\end{tabularx}
\end{table}

三者都可以理解为下面这句话：
\begin{equation}
\boxed{\text{输出}=\sum_i\text{记忆内容}_i\times\text{匹配权重}_i.}
\end{equation}
差别在于\textbf{权重怎么计算}、\textbf{Key/Value 存在哪里}，以及\textbf{怎样更新记忆}。

\subsection{MLP 和记忆型 Attention 不是同一个东西}
\begin{table}[H]
\centering\small
\renewcommand{\arraystretch}{1.4}
\begin{tabularx}{\linewidth}{@{}p{28mm}XX@{}}
\toprule
 & MLP 的参数记忆视角 & Self-Attention 的上下文记忆视角 \\
\midrule
Key / Value 来源 & 训练得到的权重 & 当前输入的 Token 表示 \\
更新方式 & 通常需要修改参数 & 可以通过加入上下文动态改变 \\
主要成本 & 模型参数与训练 & 上下文长度、计算和 KV Cache \\
数学关键 & 非线性门控、特征混合 & QK 相似度和 Softmax 加权 \\
\bottomrule
\end{tabularx}
\end{table}

所以“Transformer 的每层都是记忆系统”是一个\textbf{强有力的理解视角}，而不是声称所有层都像数据库一样一条条存储事实。

\section{复习题：确认自己真的理解了}

\begin{enumerate}[label=\textbf{题 \arabic*.}]
\item $\vect{a}=(1,0)$，$\vect{b}=(0.8,0.6)$，两者点积是多少？相似度高还是低？
\item 如果只存一条记忆 $W=\vect{v}\vect{k}^\top$ 且 $\vect{k}$ 为单位向量，那么查询 $\vect{k}$ 会得到什么？
\item 为什么 $W=\sum_i\vect{v}_i\vect{k}_i^\top$ 在查询某条记忆时可能产生误差？
\item 比较 $0.5$、$0.5^2$、$0.5^3$：越来越高次的激活为什么有助于压低弱匹配？
\item Softmax 对两个分数 $[1,0.6]$ 做指数归一化。增大 $\beta$ 后，哪条记忆的权重上升？
\item 若 $K=\begin{bmatrix}1&2\\3&4\end{bmatrix}$、$\vect{x}=\begin{bmatrix}1\\-1\end{bmatrix}$，计算 $\relu(K\vect{x})$。
\item 若 $V^\top=\begin{bmatrix}2&0\\1&3\end{bmatrix}$、$\vect{a}=\begin{bmatrix}0\\2\end{bmatrix}$，计算 $V^\top\vect{a}$，并用「按列加权求和」解释结果。
\end{enumerate}

\begin{tcolorbox}[colback=paper,colframe=linegray,boxrule=.6pt,arc=1.5mm,title={\sffamily\bfseries 简明答案},fonttitle=\sffamily\bfseries,breakable]
\begin{enumerate}[label=\arabic*.]
\item $0.8$，两单位向量方向比较接近。
\item $W\vect{k}=\vect{v}(\vect{k}^\top\vect{k})=\vect{v}$。
\item 因为其他 Key 与查询的点积可能非零，它们的 Value 也会被加进输出。
\item 它们分别是 $0.5,0.25,0.125$；高次激活能更快压低小于 $1$ 的相似度。
\item 第一条（分数为 $1$）的相对权重上升；当 $\beta$ 很大时趋近于 $1$。
\item $K\vect{x}=(-1,-1)^\top$，所以 $\relu(K\vect{x})=(0,0)^\top$。
\item $V^\top\vect{a}=(0,6)^\top$；这是两列 Value 按权重 $0$ 与 $2$ 求和。
\end{enumerate}
\end{tcolorbox}

\section*{最后记住这一条主线}
\addcontentsline{toc}{section}{最后记住这一条主线}
\begin{conclusion}[title={一张脑中的路线图}]
\textbf{点积}：度量 Query 与 Key 的匹配程度。\\
\textbf{外积}：把 Key--Value 关联写入一个矩阵。\\
\textbf{线性召回}：把匹配项和不匹配项一起混合。\\
\textbf{MLP 非线性}：抑制弱匹配，从而缓解干扰。\\
\textbf{Softmax Attention}：把检索权重集中在更相关的记忆上。\\[2mm]
\textbf{最核心的一个等式：}
\[
\boxed{W=\sum_i\vect{v}_i\vect{k}_i^\top
\quad\Longrightarrow\quad W\vect{x}=\sum_i\vect{v}_i(\vect{k}_i^\top\vect{x}).}
\]
这说明：\textbf{矩阵乘法也可以被理解为按相似度检索一组关联记忆。}
\end{conclusion}

\vspace{7mm}
{\footnotesize\color{muted}
\noindent\textbf{参考资料：}Alex Wang, \emph{You could have invented transformers from memory}.\newline
\url{https://alexwang.ai/posts/associative-memory}
\par
\noindent\textbf{说明：} 本文将联想记忆的数学分析改写为初学者可以逐步计算的教程；所用国家、向量、容量示意仅为教学案例。
}
\end{document}